\documentclass[12pt,letterpaper]{article}
\usepackage{graphicx,textcomp}
\usepackage{natbib}
\usepackage{setspace}
\usepackage{fullpage}
\usepackage{color}
\usepackage[reqno]{amsmath}
\usepackage{amsthm}
\usepackage{fancyvrb}
\usepackage{amssymb,enumerate}
\usepackage[all]{xy}
\usepackage{endnotes}
\usepackage{lscape}
\newtheorem{com}{Comment}
\usepackage{float}
\usepackage{hyperref}
\newtheorem{lem} {Lemma}
\newtheorem{prop}{Proposition}
\newtheorem{thm}{Theorem}
\newtheorem{defn}{Definition}
\newtheorem{cor}{Corollary}
\newtheorem{obs}{Observation}
\usepackage[compact]{titlesec}
\usepackage{dcolumn}
\usepackage{tikz}
\usetikzlibrary{arrows}
\usepackage{multirow}
\usepackage{xcolor}
\newcolumntype{.}{D{.}{.}{-1}}
\newcolumntype{d}[1]{D{.}{.}{#1}}
\definecolor{light-gray}{gray}{0.65}
\usepackage{url}
\usepackage{listings}
\usepackage{color}

\definecolor{codegreen}{rgb}{0,0.6,0}
\definecolor{codegray}{rgb}{0.5,0.5,0.5}
\definecolor{codepurple}{rgb}{0.58,0,0.82}
\definecolor{backcolour}{rgb}{0.95,0.95,0.92}

\lstdefinestyle{mystyle}{
	backgroundcolor=\color{backcolour},   
	commentstyle=\color{codegreen},
	keywordstyle=\color{magenta},
	numberstyle=\tiny\color{codegray},
	stringstyle=\color{codepurple},
	basicstyle=\footnotesize,
	breakatwhitespace=false,         
	breaklines=true,                 
	captionpos=b,                    
	keepspaces=true,                 
	numbers=left,                    
	numbersep=5pt,                  
	showspaces=false,                
	showstringspaces=false,
	showtabs=false,                  
	tabsize=2
}
\lstset{style=mystyle}
\newcommand{\Sref}[1]{Section~\ref{#1}}
\newtheorem{hyp}{Hypothesis}


\title{Problem Set 4}
\date{Due: November 27, 2025}
\author{Applied Stats/Quant Methods 1}


\begin{document}
	\maketitle
	\section*{Instructions}
	\begin{itemize}
		\item Please show your work! You may lose points by simply writing in the answer. If the problem requires you to execute commands in \texttt{R}, please include the code you used to get your answers. Please also include the \texttt{.R} file that contains your code. If you are not sure if work needs to be shown for a particular problem, please ask.
		\item Your homework should be submitted electronically on GitHub.
		\item This problem set is due before 23:59 on Thursday November 27, 2025. No late assignments will be accepted.
	\end{itemize}



	\vspace{.5cm}
\section*{Question 1: Economics}
\vspace{.25cm}
\noindent 	
In this question, use the \texttt{prestige} dataset in the \texttt{car} library. First, run the following commands:

\begin{verbatim}
install.packages(car)
library(car)
data(Prestige)
help(Prestige)
\end{verbatim} 


\noindent We would like to study whether individuals with higher levels of income have more prestigious jobs. Moreover, we would like to study whether professionals have more prestigious jobs than blue and white collar workers.

\newpage
\begin{enumerate}
	
	\item [(a)]
	Create a new variable \texttt{professional} by recoding the variable \texttt{type} so that professionals are coded as $1$, and blue and white collar workers are coded as $0$ (Hint: \texttt{ifelse}).
	
	\vspace{.5cm}
			\lstinputlisting[language=R, firstline=46, lastline=46]{PS04_answers_SN.R}  
	

	
	\item [(b)]
	Run a linear model with \texttt{prestige} as an outcome and \texttt{income}, \texttt{professional}, and the interaction of the two as predictors (Note: this is a continuous $\times$ dummy interaction.)
	
	\vspace{.5cm}
		\lstinputlisting[language=R, firstline=52, lastline=55]{PS04_answers_SN.R}  
		\begin{footnotesize}
		
		\begin{verbatim}
			Coefficients:
			(Intercept)               income         professional  income:professional  
			21.142259             0.003171            37.781280            -0.002326  
			
		\end{verbatim}
	\end{footnotesize}
	
	\item [(c)]
	Write the prediction equation based on the result.
	
	$$
	Y_{prestige} = 21.142  + 0.0031X_{income} + 37.781X_{professional} - 0.0023X_{income}X_{professional}	
	$$
	

	\item [(d)]
	Interpret the coefficient for \texttt{income}.
	
	On average, for a non-professional role ($X_{professional} =0 )$) with a zero income and professional status, one would expect 21.142 units of prestige.
	
	\vspace{.5cm}	
	\item [(e)]
	Interpret the coefficient for \texttt{professional}.
	
	On average, holding all other variables constant, the change between a non-professional and a professional designation (a 1 unit increase in this dummy variable)  would predict an increase of  37.781 units of prestige. 
	
	\newpage
	\item [(f)]
	What is the effect of a \$1,000 increase in income on prestige score for professional occupations? In other words, we are interested in the marginal effect of income when the variable \texttt{professional} takes the value of $1$. Calculate the change in $\hat{y}$ associated with a \$1,000 increase in income based on your answer for (c).
	
	\vspace{.5cm}
	$$
	Y_{prestige} = 21.142  + 0.0031X_{income} + 37.781X_{professional} - 0.0023X_{income}X_{professional}	
	$$
	$$
	Y_{prestige} = 62.023 - (0.0023 * 1000 * 1) =  59.723
	$$
	
	On average for professional occupations, the marginal effect of income for a \$1000 increase in salary has the effect of a 59.723 unit increase in prestige. 
	
	\item [(g)]
	What is the effect of changing one's occupations from non-professional to professional when her income is \$6,000? We are interested in the marginal effect of professional jobs when the variable \texttt{income} takes the value of $6,000$. Calculate the change in $\hat{y}$ based on your answer for (c).
	
	$$
	Y_{prestige} = 21.142  + 0.0031X_{income} + 37.781X_{professional} - 0.0023X_{income}X_{professional}	
	$$
	$$
	{delta}Y_{prestige\_prof} =37.781 -0.0023X_{income}  = 23.981
	$$
	
On average,	the professional occupation receives 23.981 more units of prestige than a non professional occupation at a sallary of \$6000.
	
\end{enumerate}

\newpage

\section*{Question 2: Political Science}
\vspace{.25cm}
\noindent 	Researchers are interested in learning the effect of all of those yard signs on voting preferences.\footnote{Donald P. Green, Jonathan	S. Krasno, Alexander Coppock, Benjamin D. Farrer,	Brandon Lenoir, Joshua N. Zingher. 2016. ``The effects of lawn signs on vote outcomes: Results from four randomized field experiments.'' Electoral Studies 41: 143-150. } Working with a campaign in Fairfax County, Virginia, 131 precincts were randomly divided into a treatment and control group. In 30 precincts, signs were posted around the precinct that read, ``For Sale: Terry McAuliffe. Don't Sellout Virgina on November 5.'' \\

Below is the result of a regression with two variables and a constant.  The dependent variable is the proportion of the vote that went to McAuliff's opponent Ken Cuccinelli. The first variable indicates whether a precinct was randomly assigned to have the sign against McAuliffe posted. The second variable indicates
a precinct that was adjacent to a precinct in the treatment group (since people in those precincts might be exposed to the signs).  \\

\vspace{.5cm}
\begin{table}[!htbp]
	\centering 
	\textbf{Impact of lawn signs on vote share}\\
	\begin{tabular}{@{\extracolsep{5pt}}lccc} 
		\\[-1.8ex] 
		\hline \\[-1.8ex]
		Precinct assigned lawn signs  (n=30)  & 0.042\\
		& (0.016) \\
		Precinct adjacent to lawn signs (n=76) & 0.042 \\
		&  (0.013) \\
		Constant  & 0.302\\
		& (0.011)
		\\
		\hline \\
	\end{tabular}\\
	\footnotesize{\textit{Notes:} $R^2$=0.094, N=131}
\end{table}

\vspace{.5cm}
\begin{enumerate}
	\item [(a)] Use the results from a linear regression to determine whether having these yard signs in a precinct affects vote share (e.g., conduct a hypothesis test with $\alpha = .05$).
	\vspace{.5cm}

	To conduct a hypothesis test, we need to determine  $ \beta_{variable} / stderr $ which will give the t score for that variable for each test.
	
 $H_{0} =>  \beta_{variable}= 0$ \hspace{1cm}	  $H_{\alpha} =>  \beta_{variable}  \ne  0$
  
  $t\_score = \beta_{variable} / stderr  =  0.042/ 0.016 = 2.625$
  
  For this case  our null hypothesis is that there is no effect of having a yard sign in a precinct which yields a $t\_score = 2.625$.    Performing a 2 sided t\_test, with  $\alpha = 0.05$,  we are testing  $ \left| t\_score \right|  $ 
  compared to the known $\alpha =0.05 /  2 => p\_value = 1.96 $ therefore,  since $ \left| t\_score \right|  > 1.96 $ we \emph{reject}  the null hypothesis.
  
  
	
	\newpage		
	\item [(b)]  Use the results to determine whether being
	next to precincts with these yard signs affects vote
	share (e.g., conduct a hypothesis test with $\alpha = .05$).
	
		\vspace{.5 cm}
	$H_{0} =>  \beta_{variable}= 0$ \hspace{1cm}	 $H_{\alpha} =>  \beta_{variable}  \ne  0$
	
	$t\_score = \beta_{variable} / stderr  =  0.042/ 0.013 = 3.231 $
	
	For this case  our null hypothesis is that there is no effect of having a yard sign in a neighboring precinct which yields a $t\_score = 3.231 $.    Performing  a 2 sided t\_test, with  $\alpha = 0.05$,  we are testing  $ \left| t\_score \right|  $ 
	compared to the known $\alpha =0.05 / 2 => p\_value = 1.96 $ therefore,  since $ \left| t\_score \right|  > 1.96 $ we \emph{reject}  the null hypothesis.
	
	\vspace{.5 cm}
	\item [(c)] Interpret the coefficient for the constant term substantively.
	\vspace{.5cm}

Given the regression table information, the Constant term is the expected proportion of the vote that went to McAuliff’s opponent Ken Cuccinelli in precincts neither with the treatment of posted signs nor adjacent to precincts with posted signs.
	
	\item [(d)] Evaluate the model fit for this regression.  What does this	tell us about the importance of yard signs versus other factors that are not modeled?
		\vspace{.5cm}
		
		$R^2$ offers a measurement of predictive power.  In this example, a $R^2$ of 0.094 tells us that there are other factors which effect a candidate's voting choice.  
		while we don't have data about the significance of the two variables tested through this experiment, it is possible for the $R^2$ to be low even when the measured variables have 
		significance within the bounds of the experiment. 
		
		\end{enumerate}  


\end{document}
